\bookchapter{第十四章：密碼管理器專題}

\chapterintro{於數字之世，密碼乃君家門之鑰匙。然鑰匙太多、太繁，奈何？密碼管理器乃君之“萬能鑰匙管家”。}

\booksection{一、何以君需密碼管理器？}

今人平均有幾十甚至上百個網絡賬號。若每賬號皆用同密碼，一旦一網站被黑，所有賬號皆危。若每賬號皆用不同之繁密碼，根本記不住。

密碼管理器解此困：君只需記一主密碼，其助君記所有其他密碼。

\booksection{二、密碼管理器之工作原理}

\begin{enumerate}
  \item \textbf{加密存儲}：所有密碼皆經高強度加密存儲，惟君能解密。
  \item \textbf{主密碼}：君設一主密碼，以解鎖密碼庫。主密碼本身不存於服務器上。
  \item \textbf{自動填充}：於網站或App登錄界面，密碼管理器自動填充賬號密碼。
  \item \textbf{密碼生成}：可生隨機、繁復之密碼，提安全性。
\end{enumerate}

\booksection{三、加密理簡介}

瞭解密碼管理器之加密原理，能幫汝判斷一個服務商是否值得信任。下面用通俗之語言解釋幾個核心概念。

\booksubsection{AES-256 加密}

AES（高級加密標準）乃美國國家標準與技術研究院（NIST）認可之對稱加密算法，爲銀行、軍方、政府廣泛使用。“256”指密鑰長度爲 256 位——理論上須超過宇宙年齡之時間才能暴力破解。主流密碼管理器（1Password、Bitwarden、KeePass 等）皆用 AES-256 加密汝之密碼庫。

\booksubsection{零知架構（Zero-Knowledge Architecture）}

“零知識”非說服務商何皆不存，而乃曰\textbf{加密和解密皆於汝之設備本地完成，服務器唯保存已加密之密文}。具體程：

\begin{enumerate}
  \item 汝於客戶端輸入主密碼。
  \item 客戶端用主密碼派生出加密密鑰（這一步於汝手機/電腦上完成）。
  \item 只有加密後之密文方上傳至服務器。
  \item 下次登錄時，服務器以密文下發，客戶端再次用主密碼解密。
\end{enumerate}

鍵點：\textbf{主密碼本身自不離開汝之設備}，服務器上亦沒有。故縱服務商爲黑，黑客拿至之亦只是一堆無法解密之密文——除非彼等能猜出汝之主密碼。

\begin{tcbinfo}{何故"零知識"意味着服務商亦觀不至汝之密碼？}
蓋因解密密鑰由汝之主密碼派生，而主密碼唯存在於汝腦中、唯於汝本地設備上使用。服務商之服務器上既沒有主密碼，亦沒有解密密鑰，只有密文。這就是何故 Bitwarden、1Password 等正規服務商皆明確聲明“我們亦無法看到汝之密碼”——非彼等不欲觀，乃技術上作不至。

\end{tcbinfo}

\booksubsection{密鑰派生函數：PBKDF2 與 Argon2}

主密碼不可直接用作加密密鑰——它須經過“密鑰派生函數”處理。此函數之作用乃\textbf{故意使計算變慢}，使黑客暴力破解之成本高至不可行。

\begin{itemize}
  \item \textbf{PBKDF2}：老牌標準，爲 iCloud 鑰匙串、Bitwarden、LastPass 等廣泛使用。通過反覆迭代（通常 10 萬次以上）哈希來拖慢破解。
  \item \textbf{Argon2}：2015 年密碼哈希競賽冠軍，比 PBKDF2 更現代。除矣"算力消耗"尚引入矣"內存消耗"——使 GPU/ASIC 等專用硬件亦難以加速破解。Bitwarden、KeePassXC 等已支持。
\end{itemize}

迭代次數越高越安全，但解鎖速度亦將變慢。正規服務商都將公開其使用之派生算法與迭代次數，便於安全社區審計。

\booksubsection{案例之鑑：LastPass 2022 年數據之洩}

2022 年 8 月，老牌密碼管理器 LastPass 公告其開發環境爲入侵；同年 11 月進一步披露，黑客拿走矣用戶之密碼庫備份（加密後之密文）。事件之關鍵教訓：

\begin{itemize}
  \item \textbf{零知識架構守住矣底線}：由於密碼庫本身乃加密之，黑客拿至密文但解不開——前提乃用戶之主密碼足夠強。
  \item \textbf{弱主密用戶最危}：黑客可離線暴力破解弱主密碼，因此使用弱主密碼或重複主密碼之用戶面臨實質風險。
  \item \textbf{元數據亦泄露}：雖密碼內容加密，但站點 URL、備註等元數據或未加密，將暴露用戶使用矣何服務。
\end{itemize}

\begin{tcbwarn}{自 LastPass 事件學至之}
第一，主密碼須長且獨特（建議 16 位以上，包含大小寫、數字、符號）；第二，優先選擇支持 Argon2id、迭代次數高之服務商；第三，對高價值賬號（銀行、郵箱）務必開啓獨立之二次驗證，勿唯依賴密碼管理器。

\end{tcbwarn}

\booksection{四、主流密管器薦}

\booksubsection{系統自帶：}

\begin{itemize}
  \item \textbf{iCloud鑰匙串}：蘋果生態用戶首選。與Safari、App無縫集成，免費，安全。弊乃跨平臺使用不便。
  \item \textbf{Google密碼管理器}：安卓用戶方便。與Chrome、安卓系統集成，免費。
  \item \textbf{華爲密碼保險箱}：鴻蒙系統內置，宜華爲用戶。
\end{itemize}

\booksubsection{第三方專業工具：}

\begin{itemize}
  \item \textbf{Bitwarden}：開源免費，功能強大，支持所有平臺。強烈薦予注重隱私之用戶。
  \item \textbf{1Password}：界面美觀，功能豐富，宜家庭與團隊。付費。
  \item \textbf{LastPass}：老牌密碼管家，免費之版其能有限。
  \item \textbf{KeePass}：完全本地存儲，開源免費。宜極客用戶。
\end{itemize}

\booksection{五、密強度真實評估}

密碼管理器普遍將顯示“密碼強度”指示條（紅/黃/綠），但此評估往往只觀長度與字符種類，並不全面。真正判斷密碼是否安全，須考慮“熵”與“是否已泄露”。

\booksubsection{熵（Entropy）：密碼強度之科學度量}

“熵”衡量密碼之不確定性，單位乃比特。熵越高，暴力破解所需嘗試次數越多。

\begin{itemize}
  \item \textbf{"123456"}：熵極低，瞬間破解。
  \item \textbf{"P@ssw0rd"}：看似複雜，但出今凡字典攻擊列表裏，熵依然低。
  \item \textbf{"correct horse battery staple"}：四個隨機單詞，熵約 44 較特，破解難度遠超"P@ssw0rd"。
  \item \textbf{16 位隨機字符（含大小寫+數+符號）}：熵約 100 較特，目前算力無法暴力破解。
\end{itemize}

結論：\textbf{“長度 + 隨機性”較“複雜符號”更要}。這亦乃密碼管理器生成之隨機長密碼真正安全之原因。

\booksubsection{HaveIBeenPwned：查汝之密碼是否已泄露}

安全研究員 Troy Hunt 維護之\href{https://haveibeenpwned.com}{haveibeenpwned.com}收錄矣歷次大型數據泄露中流出之賬號與密碼。汝可：

\begin{enumerate}
  \item 輸入郵箱，查汝之賬號是否於某次泄露中爲曝光。
  \item 用 Pwned Passwords 功能查詢某個密碼是否於已知泄露庫中（密碼管理器與瀏覽器皆已內置此查詢）。
\end{enumerate}

若一個密碼出現在泄露庫中，\textbf{無論它看起來多複雜皆須立即更換}——蓋因字典攻擊將優先嚐試此已知密碼。

\begin{tcbinfo}{k-Anonymity 查詢理}
擔心於 HaveIBeenPwned 輸入密碼將再次泄露？不必擔心。現代查詢採用“k-匿名”協議：瀏覽器唯以密碼哈希之前 5 位發給服務器，服務器返回所有以這 5 位開頭之哈希列表，本地再比對。亦就是曰，服務器永遠觀不至汝之完整密碼哈希。1Password、Bitwarden 之“已泄露密碼檢測”皆乃這樣工作之。

\end{tcbinfo}

\booksubsection{密碼管理器內置之監控功能}

\begin{itemize}
  \item \textbf{1Password Watchtower}：實時監控弱密碼、重複密碼、已泄露密碼、過期密碼、不使用 2FA 之網站，並以紅/橙/黃分級展示。
  \item \textbf{Bitwarden Vault Health}：免費版即可觀弱密、重複密；付費版可檢測已泄露密（基於 HaveIBeenPwned）。
  \item \textbf{iCloud 鑰匙串"安全宜"}：iOS 設置 密碼 安全建議，將提示汝何密碼出現在已知泄露中、何重複使用。
  \item \textbf{Google 密管器"密檢工具"}：Chrome 內置，可一鍵檢已泄露、弱、重複密。
\end{itemize}

\begin{tcbpractice}{現在遂查：汝之密碼是否安全}
\begin{enumerate}
  \item 打開 haveibeenpwned.com，輸入汝之常用郵箱，觀看出現在何泄露事件中。
  \item iPhone 用戶：設 密 安全宜，逐條處紅色警。
  \item Chrome 用戶：設 密 檢密，依示改。
  \item 1Password/Bitwarden 用戶：打開 Watchtower/Health Reports，優先處理"已泄露"與"重複"項。
\end{enumerate}

\end{tcbpractice}

\booksection{六、Passkey：無密登錄之未來}

2024-2026 年，密碼技術方經歷一場革命——\textbf{Passkey（通行密鑰）}有望徹底淘汰傳統密碼。Apple、Google、Microsoft 三大平臺聯合推動，FIDO 聯盟與萬維網聯盟（W3C）背書，這乃數十年來登錄方式最大之一次變革。

\booksubsection{何乃 Passkey？}

Passkey 基於 FIDO2 / WebAuthn 標準，用\textbf{設備上之密鑰對}替代密碼。原理可簡化爲：

\begin{itemize}
  \item 汝之手機生成一對密鑰：\textbf{公鑰}（註冊賜網站）+\textbf{私鑰}（永遠存在汝設備之安全芯片中）。
  \item 登錄時，網站發一個"挑戰"，汝之手機用私鑰簽名，再由網站用公鑰驗證。
  \item 簽名前須汝通過生物識別（Face ID / 指紋）或設備 PIN 確認身份。
  \item 整個程中，\textbf{沒有任何密碼於網絡中傳輸}。
\end{itemize}

\booksubsection{Passkey vs 傳統密}

\begin{center}
\begin{tabularx}{\linewidth}{|X|X|X|}
\hline
\textbf{維度} & \textbf{傳統密} & \textbf{Passkey} \tabularnewline
\hline
抗釣魚 & 差，釣魚網站可竊取密 & 天然抗釣魚，私鑰唯對真實域名簽名 \tabularnewline
服器泄露響 & 密碼哈希或爲破解 & 服務器唯存公鑰，泄露無影響 \tabularnewline
便性 & 需記憶/輸入/管 & Face ID/指紋一秒登錄 \tabularnewline
重放攻擊 & 或爲中間人截獲重放 & 每次簽名一性，無法重放 \tabularnewline
\hline
\end{tabularx}
\end{center}

\booksubsection{各平臺支情況}

\begin{itemize}
  \item \textbf{Apple iCloud 鑰匙串 Passkey}：iOS 16+ / macOS Ventura+ 原生支持，Passkey 可於所有登錄同一 Apple ID 之設備間同步，新設備登錄網站時無需重新創建。
  \item \textbf{Google 密管器 Passkey}：安卓 + Chrome 原生支，Passkey 存 Google 賬戶中，跨安卓設同步；2023 年起 Google 個人賬號已支用 Passkey 登錄。
  \item \textbf{Microsoft}：Windows 11 / Edge 已支持 Passkey；微軟賬號、Outlook、Office 等服務陸續開放 Passkey 登錄，並可於 Microsoft Authenticator 中管理。
  \item \textbf{1Password / Bitwarden}：已支持 Passkey 之創建、存儲、跨平臺同步，使汝不受限於單一生態。
\end{itemize}

\begin{tcbinfo}{跨設備同步乃關鍵突破}
早期 Passkey 之最大痛點乃“綁定單臺設備”——換手機遂要重新創建。2023 年起 Apple、Google、Microsoft 實現矣 Passkey 之雲端同步：iCloud 鑰匙串同步跨 Apple 設備，Google 密碼管理器同步跨安卓設備。1Password 等第三方則打通矣跨生態同步。這意味着 Passkey 已具備替代密碼之實用性。

\end{tcbinfo}

\booksubsection{何以於手機上創建與使用 Passkey}

\textbf{iOS（iPhone）徑：}

\begin{enumerate}
  \item 於支持 Passkey 之網站（如 Google、微軟、亞馬遜、eBay、PayPal）賬號安全頁選擇"添加通行密鑰 / Add passkey"。
  \item 系統彈出 Face ID / Touch ID 驗證窗口，通過後 Passkey 自動生成並保存至 iCloud 鑰匙串。
  \item 下次登錄該網站時，輸入用戶名後擇"用 iPhone 登錄"，Face ID 一刷即可。
  \item 管入：設 密 找至對應網站，可見 Passkey 條目。
\end{enumerate}

\textbf{安卓徑：}

\begin{enumerate}
  \item 於支持 Passkey 之網站賬號安全頁選擇"添加通行密鑰"。
  \item Google 密碼管理器彈窗 選擇汝之 Google 賬號 屏幕鎖驗證（指紋/圖案/PIN）。
  \item Passkey 自動保存至 Google 密碼管理器，跨同賬號安卓設備同步。
  \item 下次登錄時，擇"用 Passkey"，屏幕鎖一解即可。
  \item 管理入口：安卓設置 密碼、密鑰與自動填充 密碼管理器 google.com。
\end{enumerate}

\booksubsection{Passkey 跨設登錄}

若汝於 Windows 電腦上登錄一個只有 iPhone 上有 Passkey 之網站，可掃碼授權：

\begin{enumerate}
  \item 網站登錄頁擇"用 Passkey"。
  \item 擇"用其他設"。
  \item 電腦顯二維碼，用 iPhone 掃碼。
  \item iPhone 彈出 Face ID 驗證，確認後網站於電腦上完成登錄。
\end{enumerate}

這一過程通過藍牙+互聯網雙重加密通信，Passkey 私鑰未能離開汝之 iPhone。

\begin{tcbpractice}{現在遂試：體驗無密碼登錄}
\begin{enumerate}
  \item 登錄 Google 賬號 安全 通行密鑰 添通行密鑰，體驗 Face ID/指紋登錄。
  \item 登錄微軟賬號 安全 高級安全項 添通行密鑰。
  \item 於淘寶、京東、eBay、PayPal 等已支持 Passkey 之服務上嚐鮮。
  \item 保留傳統密碼作爲後備，逐步遷移高價值賬號至 Passkey。
\end{enumerate}

\end{tcbpractice}

\begin{tcbwarn}{過渡期宜}
Passkey 並未悉取代密——許多網站仍處於“Passkey + 密”並存階段。宜：高價值賬號優先開啓 Passkey，但\textbf{勿立即刪除傳統密碼}，留作備用；同時繼續使用密碼管理器，蓋因非所有網站皆支持 Passkey。以 Passkey 觀作“密碼管理器之強力補充”而非“立刻替換”。

\end{tcbwarn}

\booksection{七、硬件安全密鑰}

若汝對安全之要求高於普通用戶——比如持有大量加密貨幣、乃高管/記者/活動人士、或保護企業核心賬號——硬件安全密鑰乃比短信驗證碼更可靠之選擇。

\booksubsection{何乃硬件安全密鑰？}

硬件密鑰乃一個物理設備（形似 U 盤），內部安全芯片存儲汝之私鑰，須物理接觸或 NFC 才能完成驗證。代表產品：

\begin{itemize}
  \item \textbf{YubiKey}（Yubico 公司）：最主流，支 FIDO2/WebAuthn、U2F、OTP、PIV、OpenPGP，USB-C / USB-A / NFC 等多種接口。5 系列 50-70 美元。
  \item \textbf{Google Titan}：Google 自家出品，USB-C + NFC / 藍牙，約 30-40 美元。
  \item \textbf{OnlyKey}：開源替代品，支多種協議。
  \item \textbf{國產替代}：如飛天誠信 ePasskey，價更親民。
\end{itemize}

\booksubsection{何故硬件密鑰比短信驗證碼安全？}

\begin{itemize}
  \item \textbf{抗釣魚}：與 Passkey 一樣，硬件密鑰唯對真實域名響應，釣魚網站無法騙取驗證碼。
  \item \textbf{抗 SIM 劫持}：短信驗證碼可爲 SIM-swap 攻擊截獲，硬件密鑰無需手機號。
  \item \textbf{抗中間人}：FIDO2 協議每次簽名皆包含挑戰值，無法重放。
  \item \textbf{物理隔離}：私鑰永遠於硬件芯片內，惡意軟件縱攻陷汝之電腦亦讀不至。
\end{itemize}

\booksubsection{適場景}

\begin{itemize}
  \item 加密貨幣交易所賬號、財包復。
  \item 企業 SSO、遠程辦公 VPN、SSH 登錄服器。
  \item 郵箱、GitHub、社交媒體等"主賬號"——一旦爲攻破可重置所有其他賬號。
  \item 記者、律師、高管等高敏感人羣。
\end{itemize}

\booksubsection{何以配置硬件密鑰作爲二次驗證}

以 Google 賬號 + YubiKey 爲例：

\begin{enumerate}
  \item 購買 YubiKey（建議一次買兩以，一以日常用、一以備用，存放於不同地點）。
  \item 登錄 myaccount.google.com 安全 兩步驗。
  \item 選擇"添加安全密鑰" 以 YubiKey 插入 USB-C 口或貼近手機 NFC 感應區。
  \item 依示觸摸 YubiKey 金屬觸點畢註冊。
  \item 建議注冊第二以備用密鑰。
  \item 於"兩步驗證"設置裏，以硬件密鑰設爲首選驗證方式，短信驗證碼作爲備用。
\end{enumerate}

主流支持之網站：Google、GitHub、Microsoft、Apple ID、Facebook、Twitter、Dropbox、GitLab、AWS、Coinbase、Binance 等。

\begin{tcbinfo}{雙密鑰策}
硬件密鑰最大之風險乃“丟失”。建議：(1) 同時註冊兩以不同之密鑰；(2) 一以隨身攜帶（鑰匙扣），一以存放於家中保險櫃或銀行保管箱；(3) 部分網站允許以“備用恢復碼”打印出來離線保存。這樣縱主密鑰丟失，亦未能爲永久鎖於賬號外。

\end{tcbinfo}

\begin{tcbwarn}{購機渠要慎}
務必自官方店鋪或授權代理購買，勿買二手或來源不明之密鑰——理論上預置過之硬件或已爲植入後門。YubiKey 之固件於出廠後無法更新，但 Yubico 將持續推送安全公告，關注其官網即可。

\end{tcbwarn}

\booksection{八、家/團隊密共享}

家庭場景下，Netflix、家庭 Wi-Fi、路由器後臺、智能門鎖密碼須全家共用；團隊場景下，服務器、數據庫、社交媒體賬號須多人協作。\textbf{用微信發密碼、寫於共享文檔裏皆乃高危作法}——一旦聊天記錄或文檔權限泄露，密碼遂泄露矣。

\booksubsection{1Password 之“共享保險庫”（Vault）}

1Password 乃家庭與團隊共享之標杆產品。它的核心機制乃\textbf{多保險庫}：

\begin{itemize}
  \item \textbf{個人保險庫}：只有汝能觀，如銀行、郵箱等私人賬號。
  \item \textbf{共享保險庫}：家人或團隊成員皆能觀，如 Netflix、家庭 Wi-Fi、共享信用卡。
  \item \textbf{權限分級}：團隊版可設定某些成員"唯查看"或"唯查看但不編輯"。
\end{itemize}

家人或團隊成員無需共享主密碼，各自用自己之賬號登錄，但可看到爲授權之共享保險庫內容。密碼修改後所有成員實時同步。

\booksubsection{Bitwarden 之“組織”功能}

Bitwarden 之免費版支持個人密碼共享（通過“發送 Send”功能限時分享單條密碼）；付費版可創建“組織”，以若干條目共享賜家庭成員或團隊同事，權限可細分。Bitwarden 尚提供自託管選項，適合對數據主權要求高之企業。

\booksubsection{家庭場景：共享 Netflix、Wi-Fi 等密碼之安全方式}

\begin{itemize}
  \item \textbf{Netflix / 視頻網站}：用 1Password 家庭版之"共享保險庫"存登錄密碼，家庭成員各自登錄即可，無需於羣裏發。
  \item \textbf{家 Wi-Fi 密}：iOS 設備間可通過 AirDrop 直接共享 Wi-Fi 密碼賜來訪之客人，無需念出來；安卓 12+ 也有"附近分享 Wi-Fi"功能。
  \item \textbf{智能門鎖、智能家電}：儘量用廠商 App 之"家庭成員"功能爲每人單獨創建賬號，而非共享一個管理員賬號。
  \item \textbf{家用路由器後台}：管理員密碼存至 1Password 共享保險庫，訪客 Wi-Fi 單獨開一個 SSID。
\end{itemize}

\begin{tcbinfo}{1Password 家庭版之性價比}
1Password Families 約 4.99 美元/月，支持最多 5 位家庭成員，每人皆有獨立之私人保險庫 + 4 個共享保險庫。算下來人均不至 1 美元/月，乃目前家庭共享體驗最好之方案。Bitwarden 之家庭版 10 美元/年支持 6 人，性價比更高，但界面體驗略遜。

\end{tcbinfo}

\booksubsection{團隊場景：服器/數據庫密共享}

\begin{itemize}
  \item \textbf{1Password Teams / Business}：可按部門、項目組創建多保險庫，細粒度權限控制，凡訪問行爲可審計。
  \item \textbf{Bitwarden Teams / Enterprise}：開源策，可自託，企業版支 SSO 集成、目錄同步。
  \item \textbf{職責分離則}：管員密不應祇由一人掌握——設 2-3 名"應急聯繫人"，保鍵人員離職/休假時業務不中斷。
  \item \textbf{離職程}：員工離職時立即撤銷其密碼管理器訪問權限，必要時旋轉所有他接觸過之高敏感密碼。
\end{itemize}

\booksubsection{緊急訪問（Emergency Access）：萬一汝出事}

想象一個場景：汝獨自管理家庭所有重要賬號，但意外住院或不幸離世，家人無法登錄銀行、郵箱、保險賬號。\textbf{“緊急訪問”能}正是爲此設計——授權一位信任之家人/朋友，於汝無法響應時獲取汝之密碼庫。

\textbf{Bitwarden 緊急訪問設：}

\begin{enumerate}
  \item 登錄 Bitwarden Web Vault 設 緊急訪問。
  \item 添加信任聯繫人之郵箱，邀請對方接受。
  \item 設定"雲端雙因素認證"（建議 7-14 天）：自對方發起請求至汝確認之時間窗口。
  \item 若汝於等待期內不響應（拒絕或不操作），到期後對方自動獲得訪問權。
  \item 汝隨時可手動拒絕請求或撤銷緊急訪問權限。
\end{enumerate}

\textbf{1Password 緊急訪問：}

\begin{enumerate}
  \item 1Password 不直接叫"緊急訪問"，但通過"家庭/團隊組織"功能可實現——以家人加入家庭版後，將部分條目放入"共享保險庫"。
  \item 對個人版用戶，建議以主密碼告知一位高度信任之家人，或寫入密封信封存放於保險櫃。
  \item 1Password 尚提供"恢復碼"，開啓雙重認證時務必打印保存。
\end{enumerate}

\begin{tcbpractice}{現在遂設：爲家人留一以"備用鑰匙"}
\begin{enumerate}
  \item 選定一位信任之家人作爲"緊急聯繫人"。
  \item 于密碼管理器中設置緊急訪問（Bitwarden）或加入家庭共享保險庫（1Password）。
  \item 以主密碼 + 恢復碼寫於紙上，封進信封，存放於家中保險櫃或告知可信家人。
  \item 每半年回顧一，確認緊急聯繫人是否仍然適用。
\end{enumerate}

\end{tcbpractice}

\booksection{九、密管器用最佳實踐}

\begin{itemize}
  \item \textbf{設強主密碼}：主密碼乃最後防線，務必設一長且繁之密碼，並牢記。
  \item \textbf{啓雙重認證}：爲密碼管理器賬號啓雙重認證，增一層保護。
  \item \textbf{定期查弱密碼}：大多數密碼管理器有"安全審計"之能，助君覓出弱密碼與重複密碼。
  \item \textbf{勿忘主密碼}：主密碼無法找回，忘即意味着失所有密碼。議寫於紙上，置於安全之處。
  \item \textbf{緊急訪問}：設緊急聯繫人，萬一君無法訪問，信任之人可取君之密碼。
\end{itemize}

\begin{tcbwarn}{重要提醒}
重要提醒

\end{tcbwarn}

\begin{tcbpractice}{今即試之：始用密碼管理器}
\begin{enumerate}
  \item 擇一密碼管理器（薦Bitwarden或系統自帶）。
  \item 注冊賬號，設一強主密碼。
  \item 注冊賬號，設一強主密碼。
  \item 導入現有密碼（若有）或始添新密碼。
  \item 查安全審計報告，覓需改之密碼。
\end{enumerate}

\end{tcbpractice}
